\chapter[\emj{🧱}~BFS Avanzado (Grid, Multi-source y 0-1 BFS)]{\emj{🧱}~BFS Avanzado (Grid, Multi-source y 0-1 BFS)}

\begin{dsaquees}

El BFS normal explora un mapa nivel por nivel. Estas variantes son el
mismo BFS con pequeños giros para problemas muy comunes en entrevistas:

\begin{itemize}
\item Grid: el "grafo" es una cuadrícula (un laberinto o mapa). Cada
      celda se conecta con sus 4 vecinos (arriba, abajo, izquierda,
      derecha).
\item Multi-source: en vez de empezar en un punto, empiezas en VARIOS a
      la vez. Como varios incendios 🔥 propagándose simultáneamente.
\item 0-1 BFS: aristas que cuestan solo 0 o 1. Con una doble cola
      (deque) obtienes el camino más corto sin necesitar Dijkstra.
\end{itemize}

\end{dsaquees}

\begin{dsaotraforma}

Muchos problemas de "camino más corto" no vienen como grafos explícitos,
sino como matrices. Tratar cada celda como un nodo y sus 4 (u 8) vecinos
como aristas convierte el problema en un BFS estándar.

\textbf{Multi-source BFS}: metes todos los orígenes en la Queue al inicio
(distancia 0). El BFS expande todos los frentes en paralelo, dando la
distancia mínima a CUALQUIER origen. Evita correr BFS por cada fuente.

\textbf{0-1 BFS}: cuando los pesos son solo 0 o 1, una \verb|deque|
sustituye a la cola de prioridad de Dijkstra: las aristas de peso 0 van
al FRENTE y las de peso 1 al FINAL. Resultado: $O(V+E)$ en vez de
$O(E \log V)$.

\end{dsaotraforma}

\begin{dsadiagrama}
\begin{verbatim}
GRID (4 direcciones)          MULTI-SOURCE (2 incendios)

  . . . #                       Inicio: F en (0,0) y (2,3)
  . # . .        vecinos de     F . . .        dist=0 en ambos
  . . . .        (r,c):         . . . F        y crece hacia
  # . . .        (r+1,c)(r-1,c) . . . .        afuera a la vez.
                 (r,c+1)(r,c-1)

dr[] = {1,-1,0,0}   dc[] = {0,0,1,-1}
Multi-source: encola TODAS las fuentes con dist=0 al inicio.
\end{verbatim}
\end{dsadiagrama}

\begin{lstlisting}[language=C++]
BFS EN GRID
- Cada celda (r,c) es un nodo. Vecinos válidos: dentro de límites y
  no bloqueados. Usa arreglos de direcciones dr[]/dc[].

MULTI-SOURCE BFS
- Encola TODAS las fuentes con distancia 0 antes del bucle. El resto
  del BFS es idéntico.

📝 PSEUDOCÓDIGO (multi-source en grid)
──────────────────────────────────────────────────────────────

función multiSourceBFS(grid, fuentes):
    dist = matriz de -1 (no visitado)
    q = Queue
    PARA cada fuente (r,c): dist[r][c]=0; q.push((r,c))
    MIENTRAS q no vacia:
        (r,c) = q.pop()
        PARA cada direccion (dr,dc):
            nr=r+dr; nc=c+dc
            SI dentro de limites Y grid libre Y dist[nr][nc]==-1:
                dist[nr][nc] = dist[r][c]+1
                q.push((nr,nc))
    return dist

💻 CÓDIGO C++ (BFS en grid)
──────────────────────────────────────────────────────────────

#include <vector>
#include <queue>
using namespace std;

int dr[] = {1, -1, 0, 0};
int dc[] = {0, 0, 1, -1};

vector<vector<int>> bfsGrid(vector<vector<int>>& grid,
                            vector<pair<int,int>>& sources) {
    int R = grid.size(), C = grid[0].size();
    vector<vector<int>> dist(R, vector<int>(C, -1));
    queue<pair<int,int>> q;
    for (auto [r, c] : sources) { dist[r][c] = 0; q.push({r, c}); }

    while (!q.empty()) {
        auto [r, c] = q.front(); q.pop();
        for (int d = 0; d < 4; d++) {
            int nr = r + dr[d], nc = c + dc[d];
            if (nr < 0 || nr >= R || nc < 0 || nc >= C) continue;
            if (grid[nr][nc] == 1 || dist[nr][nc] != -1) continue; // muro/visto
            dist[nr][nc] = dist[r][c] + 1;
            q.push({nr, nc});
        }
    }
    return dist;
}

💻 CÓDIGO C++ (0-1 BFS con deque)
──────────────────────────────────────────────────────────────

#include <deque>
// adj[u] = lista de {vecino, peso} con peso 0 o 1
vector<int> zeroOneBFS(int V, vector<vector<pair<int,int>>>& adj, int src) {
    vector<int> dist(V, INT_MAX);
    deque<int> dq;
    dist[src] = 0; dq.push_front(src);
    while (!dq.empty()) {
        int u = dq.front(); dq.pop_front();
        for (auto [v, w] : adj[u]) {
            if (dist[u] + w < dist[v]) {
                dist[v] = dist[u] + w;
                if (w == 0) dq.push_front(v);  // peso 0 -> al frente
                else        dq.push_back(v);   // peso 1 -> al final
            }
        }
    }
    return dist;
}

⏱️ COMPLEJIDAD (Big-O)
──────────────────────────────────────────────────────────────

  Variante        │ Tiempo      │ Nota
  ────────────────┼─────────────┼────────────────────────
  BFS en grid     │ O(R * C)    │ R filas, C columnas
  Multi-source    │ O(R * C)    │ igual, varios orígenes
  0-1 BFS         │ O(V + E)    │ reemplaza a Dijkstra
\end{lstlisting}

\begin{dsaentrevistas}

✅ "Rotting Oranges" (LeetCode 994): multi-source BFS puro. Todas las
   naranjas podridas son fuentes; el resultado es el tiempo hasta podrir
   todo.

💡 "01 Matrix" (LeetCode 542): distancia de cada celda al 0 más cercano
   = multi-source BFS desde todos los ceros a la vez.

⚠️ 8 direcciones vs 4: lee bien el enunciado. Si permiten diagonales,
   agrega las 4 direcciones diagonales a dr[]/dc[].

🔑 0-1 BFS gana a Dijkstra cuando los pesos son binarios: mismo
   resultado, más rápido y sin heap. Buen detalle para destacar.

\end{dsaentrevistas}

\begin{dsaresumen}

• Grid = grafo implícito: cada celda es un nodo, 4 u 8 vecinos.
• Usa dr[]/dc[] para recorrer direcciones limpiamente.
• Multi-source: encola TODAS las fuentes con dist 0 al inicio.
• 0-1 BFS: deque, peso 0 al frente, peso 1 al final; O(V+E).
• Clásicos: Rotting Oranges, 01 Matrix, laberintos.
════════════════════════════════════════════════════════════════

\end{dsaresumen}
